% !PRINTED 168
\weberpage{168}
%<<CONT
wherein $\lambda_1,\ldots\lambda_n$ are constants, over which we dispose in such a way that the coefficient of $y^r$ in $F_0$ does not vanish, i.e.\ that $F_0(1,\lambda_1,\ldots\lambda_n)$ becomes different from zero. Then the system $\mathfrak S$, like every form of the variables (10) generally, passes into a system of forms of the variables $y,x_1,\ldots x_n$, and conversely every form that depends on these variables can also be represented as a form of the linear combinations (10).

Any form $F$ of the system $\mathfrak S$ is now ordered according to powers of $y$, and then the division by $F_0$ is carried out with respect to $y$, whereby
\weq{11}{F = a_0F_0 + \Phi}
may result, so that $a_0$ is the quotient and $\Phi$ the remainder of the division. $a_0$ and $\Phi$ are \gesperrt{integral functions} of the variables $x, y$, because the coefficient of the highest power of $y$ in the divisor $F_0$ is constant. $\Phi$ does not exceed the degree $r-1$ with respect to $y$. Now if $F$ runs through the system $\mathfrak S$, the totality of the functions $\Phi$ defined by (11) forms a system $\mathfrak S_{r-1}$, of which we assume the representability by a basis as already proved. We can therefore set:
\weq{12}{\Phi = A_1\Phi_1 + \cdots + A_\mu\Phi_\mu,}
so that $\Phi_1,\ldots\Phi_\mu$ belong to the system $\mathfrak S_{r-1}$, i.e.\ in such a way that in the system $\mathfrak S$ the forms
\weq{13}{F_1 = a_1F_0+\Phi_1,\ \ldots,\ F_\mu = a_\mu F_0+\Phi_\mu}
can be determined. If we therefore set
\weq{14}{A_0 = a_0 - a_1A_1 - \cdots - a_\mu A_\mu,}
then from (11), (12) and (13) it follows:
\weq{15}{F = A_0F_0+A_1F_1+\cdots+A_\mu F_\mu,}
whereby Theorem I.\ is proved in general.

\begin{center}
\S. 42.\\
\gesperrt{Finiteness of the system of invariants of a finite linear substitution group.}
\end{center}

The proof of Theorem I.\ given in the preceding section is in itself subject to no exceptions whatever. But we lose nothing essential of its generality if we exclude once and for all identically vanishing forms. If
%CONT>>
